\documentclass[10pt,a4paper]{amsart}
\usepackage[margin=19mm]{geometry}
\usepackage{amsmath,amssymb,mathtools}
\usepackage[scr=rsfs,cal=euler]{mathalfa}
\usepackage{xfrac}
\usepackage[unicode,hidelinks,bookmarks=false]{hyperref}
\newcommand{\ie}{i.e.\ }
\newcommand{\cf}{cf.\ }
\newcommand{\resp}{respectively\ }
\newcommand{\Subsection}{\subsection}
\providecommand{\nobreakdash}{\nobreak\hbox{-}}
\setlength{\emergencystretch}{2em}
\multlinegap0pt
\usepackage{amssymb}
\usepackage{xcolor}
\newtheorem{lemma}{Lemma}
\newtheorem{proposition}{Proposition}
\renewcommand{\epsilon}{\varepsilon}
\title{Existence of multi-peak solutions for singularly perturbed problems in the zero-mass case}
\author{L\'eo Demuth}
\date{Source: arXiv:2609.03135v1 (CC BY 4.0)}
\begin{document}
\maketitle

\noindent\emph{Source and licence.} Existence of multi-peak solutions for singularly perturbed problems in the zero-mass case. Version arXiv:2609.03135v1 (CC BY 4.0), distributed by arXiv under the Creative Commons Attribution 4.0 International licence, http://creativecommons.org/licenses/by/4.0/, as recorded in the arXiv licence metadata for this exact version and shown on the abstract page https://arxiv.org/abs/2609.03135v1. Full licence notice: \texttt{licenses/arxiv-2609.03135-CC-BY-4.0.txt}.\par\smallskip

\noindent\textit{Editorial scope.} Counted unit: the article's proposition Uniformity (source0.tex lines 847--855). The article's complete original proof of it is delivered with it (source0.tex lines 857--988, one \texttt{proof} environment of 132 lines). No mathematics is written, repaired or re-derived: the statement and the proof are the article's own lines, in the article's own notation, and no retained line is substituted or re-flowed.  Retained uncounted as the source-local context the proof needs: lines 410--417; lines 772--777; lines 824--837.  Nothing was omitted from any retained span, so the package carries no omission record.  No retained line contains a citation, so the wrapper carries no bibliography and no bibliography entry.  No retained line contains a figure environment, an image-inclusion command or drawing code, so no image is delivered and the package declares no asset dependency.  Editorial formatting only, in addition: an AMS \texttt{amsart} preamble that restates the article's own declarations and macros used by the retained lines, namely the environments \texttt{lemma}, \texttt{proposition} on separate counters, together with the standard packages those lines need.  The retained environments print in this document as Lemma 1, Proposition 1 in order of appearance, whereas the article prints the same items with the numbering of its own full document.  The complete native source of the article as distributed in the arXiv e-print for this exact version is bundled unchanged as \texttt{sources/209268/source0.tex}.
\begin{lemma}\label{non-degeneracy of the solution}
Consider $P\in \mathbb{R}^n$. The solutions in $D^{1,2}(\mathbb{R}^n)=\left\{u \in L^{2^*}(\mathbb{R}^n) | \nabla u \in L^2(\mathbb{R}^n) \right \}$ of the linearized problem 
\begin{equation*}\label{linear} 
   -\Delta Z=pU_P^{p-1}Z-qQ(P)U_P^{q-1}Z\quad  \text{in} \quad \mathbb{R}^n.
\end{equation*} 
are linear combinations of 
$\partial U_P\over\partial x_j$,  $j=1,\dots,n$. 
\end{lemma}

\begin{lemma}\label{ContinuityP}
$\\$

Consider a compact set $E$. Then $P\in E \rightarrow  L_{P} \in L(X)$ is continuous where $L_P\phi = \phi -i^*((f'(U_P)-g'(U_P)Q(P))\phi)$. 
$$\\$$
\end{lemma}

\begin{lemma}\label{Projection prop}

Consider two compact sets $E_1$ and $E_2$  of $~\mathbb{R}^{n}$ such that $d(E_1,E_2)>0$ and set $E:=E_1\times E_2$. 
Denote $\tau_{\epsilon,P}:\phi \in X \mapsto \phi(. +\frac{P}{\epsilon}) $. One has
\begin{itemize}
    \item There exists $\epsilon_0 >0$, There exists $C>0$ such that $\forall 0<\epsilon <\epsilon_0,\forall P\in E, ||\Pi_{\epsilon,P}||_X\leq C$. 
    
    \item There exists $\epsilon_0$ such that for all $\epsilon <\epsilon_0$, $P \mapsto \Pi_{\epsilon,P}$ is continuous.
    \item Consider sequences $(\epsilon_m)\in (\mathbb{R}_+^*)^\mathbb{N}$ and $P_m \in E^\mathbb{N}$ such that $\epsilon_m \rightarrow 0$ and $P_m \rightarrow (A_1,A_2)$. Then, $\tau_{\epsilon_m,P_{m,k}}\Pi_{\epsilon_m,P_m}\tau_{\epsilon_m,P_{m,k}}^{-1} \eta \rightarrow \Pi_{\infty,A_k}\eta$ in $X$ for all  $\eta \in C_c^\infty(\mathbb{R}^n)$ where $\Pi_{\infty,A_k}$ is the orthogonal projection on $ K_{\infty,A_k}^{\perp}$ where  $K_{\infty,A_k}=Span\{\frac{\partial U_{A_k}}{\partial x_j},j\in\{1,...,n\}\}$.
    \end{itemize}



\end{lemma}

\begin{proposition}\label{Uniformity}

Consider two compact sets $E_1$ and $E_2$  of $~\mathbb{R}^{n}$ such that $d(E_1,E_2)>0$ and set $E:=E_1\times E_2$. Then, there exists $\epsilon_0$ such that  there exists $C>0$ such that


$$\forall P\in E, 0<\epsilon <\epsilon_0,\forall\phi \in K_{\epsilon,P}^\perp, ||\tilde{L}_{\epsilon,P} \phi||_X \geq C||\phi||_X$$

$$\\$$
\end{proposition}

\begin{proof} 

    


$\\$$\\$

Assume by contradiction the existence of sequences $(\epsilon_m) \in (\mathbb{R}_+^*)^{\mathbb{N}}, (P_m) \in E^{\mathbb{N}}, (\phi_m) \in (K_{\epsilon_m,P_m}^\perp)^{\mathbb{N}}$ such that $\epsilon_m \rightarrow 0$, $||\phi_m||_X=1$ and $||\tilde{L}_{\epsilon_m,P_m}\phi_m||_X \leq \frac{1}{m}$. By compactness, we can suppose that $P_{m,k} \rightarrow A_k$ for $k\in \{1,2\}$. Fix $k_0\in \{1,2\}$. 
$\\$
Let $\tilde{\phi}_{m,k_0} (x)=\phi_m(x+\frac{P_{m,k_0} }{\epsilon_m})$. Thus

$$||\tilde{\phi}_{m,k_0} ||_X= ||\phi_{m}||_X=1.$$


Hence, by reflexivity of $X$, there exists $\psi_{k_0} \in X$ such that $\tilde{\phi}_{m,{k_0}} \rightharpoonup \psi_{k_0} \in X$. We want to show that



$$\psi_{k_0} = i^*((f'(U_{A_{k_0}})-Q(A_{k_0})g'(U_{A_{k_0}}))\psi_{k_0}). $$
Fix $\eta \in C_c^{\infty}(\mathbb{R}^n)$. Denote $\tau$ the isometry 
\begin{equation*}
\begin{cases}

X \rightarrow X \\
 \phi \mapsto \phi(.+\frac{P_{m,{k_0}}}{\epsilon_m})
\end{cases}.
\end{equation*}
We claim that
because $\tau$ is an isometry
$$\langle\tau\tilde{L}_{\epsilon_m,P_m}\phi_m,\eta\rangle=o(1).$$
Thus, because $ \Pi_{\epsilon_m,P_m}\phi_{m}=\phi_m$,
$$\left\langle\left(\tilde{\phi}_{m,{k_0}} -\tau\Pi_{\epsilon_m,P_m}\sum i^*((f'(\ U_{\epsilon,P_{m,k}})-Q(P_{m,k})g'( U_{\epsilon,P_{m,k}}){\phi}_{m})\right)|\eta\right\rangle  =o(1).$$
But, $\langle \tilde{\phi}_{m,{k_0}}|\eta\rangle\rightarrow \langle\psi_{k_0}|\eta\rangle$ by weak convergence. In addition,


$$\left\langle\tau\Pi_{\epsilon_m,P_m}\sum i^*((f'(\ U_{\epsilon,P_{m,k}})-Q(P_{m,k})g'( U_{\epsilon,P_{m,k}}){\phi}_{m})|\eta\right\rangle$$$$=\left\langle\sum i^*((f'(\ U_{\epsilon,P_{m,k}})-Q(P_{m,k})g'( U_{\epsilon,P_{m,k}}){\phi}_{m})|\Pi_{\epsilon_m,P_m}\tau^{-1}\eta\right\rangle$$
$$=\left\langle\sum (f'(\ U_{\epsilon,P_{m,k}})-Q(P_{m,k})g'( U_{\epsilon,P_{m,k}}){\phi}_{m})|\Pi_{\epsilon_m,P_m}\tau^{-1}\eta\right\rangle_2$$
( by definition of $i^*$)
$$=\left\langle \sum \left(f'( U_{P_{m,k}})\left(x-\frac{P_{m,k}-P_{m,{k_0}}}{\epsilon_m}\right)-Q(P_{m,k})g'( U_{P_{m,k}})\left(x-\frac{P_{m,k}-P_{m,{k_0}}}{\epsilon_m}\right)\right)\tilde{\phi}_{m,{k_0}}|\tau\Pi_{\epsilon_m,P_m}\tau^{-1}\eta\ \right\rangle_2$$





$\\$


Furthermore,
$$||f'( U_{P_{m,{k_0}}})-Q(P_{m,{k_0}})g'( U_{P_{m,{k_0}}})\tilde{\phi}_{m,{k_0}}
-(f'(U_{A_{k_0}})-Q(A_{k_0})g'(U_{A_{k_0}}))\psi_{k_0}||_\frac{2n}{n+2} \rightarrow 0.$$

 

Indeed,
$$||(f'( U_{P_{m,{k_0}}})-Q( P_{m,{k_0}})g'( U_{P_{m,{k_0}}}))\tilde{\phi}_{m,{k_0}}
-(f'(U_{A_{k_0}})-Q(A_{k_0})g'(U_{A_{k_0}}))\psi_{k_0}||_\frac{2n}{n+2} $$
$$
\leq ||(f'(\ U_{P_{m,{k_0}}})-f'(U_{A_{k_0}})-[Q(P_{m,{k_0}})\ g'( U_{P_{m,{k_0}}})-Q(A_{k_0})g'(U_{A_{k_0}})])\tilde{\phi}_{m,{k_0}}
||_\frac{2n}{n+2} $$ 

 $$+||(f'(U_{A_{k_0}})-Q(A_{k_0})g'(U_{A_{k_0}}))(\tilde{\phi}_{m,{k_0}}-\psi_{k_0})||_{\frac{2n}{n+2}}.$$

But $h \in X \mapsto (f'(U_{A_{k_0}})-Q(A_{k_0})g'(U_{A_{k_0}}))h) \in L^{\frac{2n}{n+2}}$ is a compact operator. Thus, $$||(f'(U_{A_{k_0}})-Q(A_{k_0})g'(U_{A_{k_0}}))(\tilde{\phi}_{m,{k_0}}-\psi_{k_0})||_{\frac{2n}{n+2}} \rightarrow 0.$$

Now, let $D_m =f'(\ U_{P_{m,{k_0}}})-f'(U_{A_{k_0}})-[Q(P_{m,{k_0}})\ g'( U_{P_{m,{k_0}}})-Q(A_{k_0})g'(U_{A_{k_0}})])$.
 One has
$$||D_m \tilde{\phi}_{m,{k_0}}||_{ \frac{2n}{n+2}} \rightarrow 0$$
 Indeed,
by the inequality of Hölder with $\frac{1}{\frac{2n}{n+2}}=\frac{1}{2}+\frac{1}{n}$,

$$||D_m\tilde{\phi}_{m,{k_0}}||_{\frac{2n}{n+2}} \leq C||\tilde{\phi}_{m,{k_0}}||_X ||D_{m}||_n $$
But  $||\tilde{\phi}_{m,{k_0}}||=1$. Also, recall that $U_P(x)=Q(P)^aU(Q(P)^bx)$,$Q(P) \in [\nu,M]$ and $U$ is decreasing. Thus $|D_m|\leq C[|U(\nu^b.)^{p-1}|+|U(\nu^b.)^{q-1}| $ and $$|||U(\nu^b.)^{p-1}|+|U(\nu^b.)^{q-1}|~||_n\leq C'||U^{q-1}||_n<+\infty.$$ In addition, by continuity of $Q$ and $U$, $D_m \rightarrow 0$ almost everywhere. We conclude by the dominated convergence theorem that $$||D_{m}||_n\rightarrow 0.$$ Then, 
$$||D_m \tilde{\phi}_{m,{k_0}}||_{ \frac{2n}{n+2}} \rightarrow 0.$$

$\\$


In addition thanks to the lemma \ref{Projection prop}, $$||\tau\Pi_{\epsilon_m,P_m}\tau^{-1}\eta -\Pi_{\infty,A_{k_0}}\eta||_{2^*}\rightarrow0.$$

Thus, by the inequality of Holder, 
$$\left\langle \tau \Pi_{\epsilon_m,P_m}i^*((f'(\ U_{\epsilon,P_{m,{k_0}}})-Q(P_{m,{k_0}})g'( U_{\epsilon,P_{m,{k_0}}}){\phi}_{m}))|\eta\right\rangle$$$$\rightarrow \langle\Pi_{\infty,A_{k_0}} i^*((f'(U_{A_{k_0}})-Q(A_{k_0})g'(U_{A_{k_0}}))\psi_{k_0})|\eta \rangle $$
Finally, when $k\ne k_0$,
$$\left|\left\langle  \left(f'( U_{P_{m,k}})\left(x-\frac{P_{m,k}-P_{m,{k_0}}}{\epsilon_m}\right)-Q(P_{m,k})g'( U_{P_{m,k}})\left(x-\frac{P_{m,k}-P_{m,{k_0}}}{\epsilon_m}\right)\right)\tilde{\phi}_{m,{k_0}}|\tau\Pi_{\epsilon_m,P_m}\tau^{-1}\eta\right\rangle_2\right|$$
$$=\left|\left\langle  \left(f'( U_{P_{m,{k_0}}})\left(x-\frac{P_{m,k}-P_{m,{k_0}}}{\epsilon_m}\right)-Q(P_{m,{k}})g'( U_{P_{m,k}})\left(x-\frac{P_{m,k}-P_{m,{k_0}}}{\epsilon_m}\right)\right)\tilde{\phi}_{m,{k_0}}|\Pi_{\infty,A_{k_0}}\eta\right\rangle_2\right|+o(1)$$
$$\leq ||\tilde{\phi}_{m,{k_0}}||_2\left|\left| \left(f'( U_{P_{m,k}})\left(x-\frac{P_{m,k}-P_{m,{k_0}}}{\epsilon_m}\right)-Q(P_{m,k})g'( U_{P_{m,k}})\left(x-\frac{P_{m,k}-P_{m,{k_0}}}{\epsilon_m}\right)\right)\tilde{\phi}_{m,{k_0}})\Pi_{\infty,A_{k_0}}\eta\right|\right|_2 $$$$+o(1) $$
But, $||\Pi_{\infty,A_{k_0}}\eta||_2<+\infty $. Then, by the dominated convergence Theorem,
$$\left|\left| \left(f'( U_{P_{m,{k}}})\left(x-\frac{P_{m,k}-P_{m,{k_0}}}{\epsilon_m}\right)-Q(P_{m,k})g'( U_{P_{m,k}})\left(x-\frac{P_{m,k}-P_{m,{k_0}}}{\epsilon_m}\right)\right)\Pi_{\infty,A_{k_0}}\eta\right|\right|_2 \rightarrow 0.$$

$$\left|\left\langle  \left(f'( U_{P_{m,k}})\left(x-\frac{P_{m,k}-P_{m,{k_0}}}{\epsilon_m}\right)-Q(P_{m,{k_0}})g'( U_{P_{m,k}})\left(x-\frac{P_{m,{k}}-P_{m,{k_0}}}{\epsilon_m}\right)\right)\tilde{\phi}_{m,{k_0}}|\tau\Pi_{\epsilon_m,P_m}\tau^{-1}\eta\right\rangle_2\right|\rightarrow0$$

Conclusion, 
$$\left\langle\left(\tilde{\phi}_{m,{k_0}} -\tau\Pi_{\epsilon_m,P_m}\sum i^*((f'(\ U_{\epsilon,P_{m,k}})-Q(P_{m,k})g'( U_{\epsilon,P_{m,k}}){\phi}_{m})\right)|\eta\right\rangle \rightarrow $$$$\langle \psi_{k_0} -\Pi_{\infty,A_{k_0}} i^*((f'(U_{A_{k_0}})-Q(A_{k_0})g'(U_{A_{k_0}}))\psi_{k_0})|\eta\rangle$$
In conclusion, by uniqueness of the limit and because it is true for all $\eta$,
$$\psi_{k_0} = \Pi_{\infty,A_{k_0}}i^*((f'(U_{A_{k_0}})-Q(A_{k_0})g'(U_{A_{k_0}}))\psi_{k_0}) $$
Furthermore, since $\phi_m \in K_{\epsilon_m,P_m}^\perp$
$$\left\langle \phi_{m}\big| \frac{\partial U_{\epsilon_m,P_{m,k_0}}}{\partial x_j}\right\rangle=0$$
$$\left\langle \tilde{\phi}_{m,{k_0}}\big| \frac{\partial U_{P_{m,k_0}}}{\partial x_j}\right\rangle=0$$
But 

$$\left|\left|\frac{\partial U_{P_{m,{k_0}}}}{\partial x_j}-\frac{\partial U_{A_{k_0}}}{\partial x_j}\right|\right|_{D^{1,2}(\mathbb{R}^n)}\rightarrow 0$$ 
and 
$$\left\langle \tilde{\phi}_{m,{k_0}}| \frac{\partial U_{A_{k_0}}}{\partial x_j}\right\rangle \rightarrow \left\langle \psi_{k_0}|\frac{\partial U_{A_{k_0}}}{\partial x_j}\right  \rangle$$
Thus, taking the limit

$$\psi_{k_0} \in K_{\infty,A_{k_0}}^\perp$$


where $K_{\infty,A_{k_0}}^\perp=\{v \in X, \langle v|\frac{\partial U_{A_{k_0}}}{\partial x_j}\rangle=0\}$. Set, $L_{\infty,A_{k_0}}(\phi)=\phi - i^*((f'(U_{A_{k_0}})-Q(A_{k_0})g'(U_{A_{k_0}}))\phi)$.
Then,

$$\Pi_{\infty,A_{k_0}}L_{\infty,A_{k_0}}\psi_{k_0} =0$$
But, thanks to lemma \ref{non-degeneracy of the solution}, $$\left\langle L_{\infty,A_{k_0}}\psi_{k_0}|\frac{\partial U_{A_{k_0}}}{\partial x_j}\right\rangle = \left\langle \psi_{k_0}|L_{\infty,A_{k_0}}\frac{\partial U_{A_{k_0}}}{\partial x_j}\right\rangle=0 $$
Then, $$L_{\infty,A_{k_0}}\psi_{k_0} =0$$
thanks to the lemma \ref{non-degeneracy of the solution}, $\psi_{k_0}=0$.  


Now,
$$ \frac{1}{m} \geq ||\tilde{L}_{\epsilon_m,P_m} \phi_m||_X \geq ||\phi_m||- ||\Pi_{\epsilon_m,P_m}\sum i^*((f'(\ U_{\epsilon_m,P_{m,k}})-Q(P_{m,k})g'( U_{\epsilon_m,P_{m,k}}){\phi}_{m})||_X $$ $$= 1 - ||\Pi_{\epsilon_m,P_m}\sum i^*((f'(\ U_{\epsilon_m,P_{m,k}})-Q(P_{m,k})g'( U_{\epsilon,P_{m,k}}))\phi_{m})||_X.$$
 Thanks to lemma \ref{Projection prop}, it suffices to show that $||\sum i^*((f'(\ U_{\epsilon,P_{m,k}})-Q(P_{m,k})g'( U_{\epsilon,P_{m,k}}))\phi_{m})||_X \rightarrow 0$ to conclude to a contradiction.

But with $\beta=\frac{ns}{n+2s}$ or $\beta=\frac{2n}{n+2}$ and $k_0 \in \{1,2\}$
$$ ||(f'(\ U_{\epsilon_m,P_{m,{k_0}}})-Q(P_{m,{k_0}})g'( U_{\epsilon_m,P_{m,{k_0}}}))\phi_{m})||_\beta$$$$=||(f'(\ U_{P_{m,{k_0}}})-Q(P_{m,{k_0}})g'( U_{P_{m,{k_0}}}))\tilde{\phi}_{m,{k_0}})||_\beta=
||(f'( U_{A_{k_0}})-Q(A_{k_0})g'( U_{A_{k_0}}))\tilde{\phi}_{m,{k_0}}||_\beta+o(1)$$

because, $P \rightarrow f'( U_{P})-Q(P)g'( U_{P})$ is continuous according to the lemma \ref{ContinuityP} and $$||\tilde{\phi}_{m,{k_0}}||_X=1.$$
But 
$h \mapsto (f'( U_{A_{k_0}})-Q(A_{k_0})g'( U_{A_{k_0}}))h $ is compact. Thus,
$$ ||(f'(\ U_{\epsilon_m,P_{m,{k_0}}})-Q(P_{m,{k_0}})g'( U_{\epsilon_m,P_{m,{k_0}}}))\phi_{m}) ||_\beta \rightarrow 0.$$

$\\$

\end{proof}

\end{document}
